\section{Extended error analysis}
This appendix dissects the largest residual errors of each tool and identifies
their causes; every claim traces to a committed artifact in \texttt{results/}.

\subsection{Tool 10: the CD15 outlier}
Across every model in the break-attempt ladder, CD15 is the worst-predicted
V1 gene (v3: $\rho = 0.155$ against a cross-gene mean near $0.51$; v4:
$0.278$). Three candidate explanations, in decreasing order of support:
(i) \emph{assay-window mismatch} --- CD15 guide activities in V1 are
compressed toward the middle of the percent-rank scale, so the rank
information available to any regressor is intrinsically weaker;
(ii) \emph{sequence-context shift} --- the CD15 target contexts have GC
content skewed relative to the RES training pool, placing much of the CD15
test mass outside the feature region where the training fit is tight;
(iii) \emph{gene-specific biology} --- chromatin accessibility at the CD15
locus in the mouse line used for V1 is not represented in the human RES
selections. Explanations (i) and (ii) are checkable from the data we ship;
(iii) is a hypothesis. We do not exclude CD15 from headline numbers.

\subsection{Tool 10: the cross-species transfer penalty is real}
Before the contamination fix, our CNN appeared to transfer from human RES to
mouse V1 with no penalty ($0.665 \to 0.781$); after removing the 1{,}836
overlapping guides the same model drops to $0.358$ on mouse --- a $46\%$
relative penalty. The ridge and GNN baselines are flat at $\approx 0.25$ on
both sides, meaning they learn less human signal to lose. The clean ordering
(deep model learns more human-specific signal and loses more of it in
transfer) is itself a finding about what these models capture: a large share
of learnable on-target signal is species- or assay-specific, not universal
sequence determinism.

\subsection{Tool 7: microhomology-deletion share resists sequence-only models}
Our tuned GBT reaches Spearman $0.210$ on MH-deletion share; enumerating
microhomology pairs directly and scoring them by length and GC gives
$0.01$--$0.11$. The gap between the learned model and the mechanistic
enumerator suggests the enumerator's assumption (deletion probability
factorizes over MH length/composition) is wrong in the inDelphi data, where
deletion spectra show strong position-dependent context effects that a
pairwise enumerator cannot express. The negative result is preserved in
\texttt{results/mhdel\_improve.json} with both variants.

\subsection{Tool 7: why $+1$ insertion is the easiest target}
The $+1$-insertion frequency achieves $\rho = 0.60$ because it is dominated
by a single mechanistic variable --- the identity of the nucleotides flanking
the cut site, which template the single-base insertion. Position-specific
one-hots capture it directly; frameshift and MH-deletion outcomes depend on
longer-range context that our feature families only partially express.

\subsection{Tool 6: where the duplex model over-flags}
RNA-duplex scoring treats guide:guide pairing as concentration-independent.
In real multiplexed transfections, guides are expressed from separate
transcripts at unequal levels; a thermodynamically plausible duplex between a
highly expressed and a weakly expressed guide may never form at an effective
rate. Our screener therefore over-flags relative to any kinetic ground truth.
We mitigate by reporting both an energy-ranked list and the calibrated binary
flag, and by quantifying the flag-rate enrichment against cross-gene and
random-spacer controls rather than claiming absolute interference rates.

\subsection{Tool 8: validation breadth}
The rearrangement-risk flagger was initially validated on the three
Kosicki~2018 loci only. The RefSeq panel
(\texttt{results/riskflag\_panel.json}) broadens this to 76 mRNA
accessions (43 human, 33 mouse) with 334 candidate NGG cut sites scored:
$61.1\%$ flag LOW, $31.7\%$ MODERATE, $7.2\%$ HIGH. The highest
mean-risk loci (VHL $0.655$, ALB $0.455$, BRAF $0.397$, PCSK9 $0.383$,
MAPT $0.357$) are repeat-rich or low-complexity genes, consistent with the
mechanistic terms. Large-deletion propensity is known to be
locus-dependent (chromatin, replication timing), so the score remains a
within-locus ranking device, not a calibrated cross-genome probability,
until long-read repair-outcome sets are integrated.

\subsection{Tool 9: the biosecurity wrapper's blind spot}
Because the Common Mechanism's multi-GB databases cannot run on free-tier
hardware, Tool 9's local pre-screen is a homology-and-motif fast filter whose
output is advisory; the authoritative decision is delegated to a full Common
Mechanism run at deployment. The wrapper's value is pipeline ergonomics
(one command, one audit record), not a new screening capability --- stated
plainly to avoid overclaiming.
